\DocumentMetadata{
  pdfversion=2.0,
  pdfstandard=ua-2,
  lang=en-US,
  tagging=on,
  tagging-setup={math/setup=mathml-SE,tabsorder=structure}
}
\documentclass[11pt,letterpaper]{article}
\usepackage[margin=0.68in,includefoot]{geometry}
\usepackage{newcomputermodern}
\usepackage{amsmath,amscd,mathtools}
\providecommand{\square}{\mdlgwhtsquare}
\usepackage[hidelinks]{hyperref}
\hypersetup{
  pdftitle={Analysis PhD Exam, May 1994},
  pdfauthor={University of Florida Department of Mathematics},
  pdfsubject={Graduate mathematics examination},
  pdfdisplaydoctitle=true,
  bookmarksopen=true
}
\setcounter{secnumdepth}{0}
\setlength{\parindent}{0pt}
\setlength{\parskip}{0.28em}
\setlength{\emergencystretch}{3em}
\setlength{\leftmargini}{2.15em}
\setlength{\leftmarginii}{2.65em}
\setlength{\leftmarginiii}{3em}
\renewcommand{\labelenumi}{\textbf{\theenumi.}}
\pagestyle{empty}
\raggedbottom
\allowdisplaybreaks
\newenvironment{examproblems}[1]{%
  \begin{enumerate}%
  \setcounter{enumi}{#1}%
  \setlength{\topsep}{0.35em}%
  \setlength{\partopsep}{0pt}%
  \setlength{\itemsep}{0.48em}%
  \setlength{\parsep}{0.18em}%
}{\end{enumerate}}
\newenvironment{examsubparts}{%
  \begin{enumerate}%
  \setlength{\topsep}{0.18em}%
  \setlength{\partopsep}{0pt}%
  \setlength{\itemsep}{0.16em}%
  \setlength{\parsep}{0.1em}%
}{\end{enumerate}}
\newenvironment{examinstructions}{%
  \par\begingroup\itshape
}{%
  \par\endgroup\vspace{0.18em}
}
\makeatletter
\renewcommand\section{\@startsection{section}{1}{0pt}%
  {-1.25ex plus -.25ex minus -.1ex}%
  {0.55ex plus .1ex}%
  {\normalfont\large\bfseries}}
\renewcommand{\@maketitle}{%
  \newpage\null\vskip -1.2em%
  {\footnotesize\bfseries
    \href{https://www.ufl.edu/}{University of Florida}, \href{https://math.ufl.edu/}{Department of Mathematics}\par}%
  \vskip 0.18em%
  \tagstructbegin{tag=Title}%
    {\LARGE\bfseries\href{https://gma.math.ufl.edu/exams/analysis-phd/}{Analysis PhD Exam}, May 1994\par}%
  \tagstructend%
  \vskip 0.35em%
  \tagmcbegin{artifact}%
    \hrule height 1pt%
  \tagmcend%
  \vskip 0.55em%
}
\makeatother
\title{Analysis PhD Exam, May 1994}
\author{}
\date{}
\begin{document}
\maketitle
\thispagestyle{empty}
\section{I.}
\begin{examinstructions}
State the following theorems:
\end{examinstructions}
\begin{examproblems}{0}
\item
The Uniform Boundedness theorem
\item
The Hahn–Banach theorem
\item
The Fatou Lemma.
\item
The Lebesgue dominated convergence theorem
\item
The Vitali theorem
\item
The Radon–Nikodym theorem
\item
The Egorov theorem
\item
The Fubini theorem
\item
The dual of \(L^p\) \((1 \leq p < \infty)\)
\item
The Lusin theorem
\item
The Riesz representation theorem in locally compact spaces
\item
The Lebesgue decomposition theorem
\end{examproblems}
\section{II}
\begin{examinstructions}
Prove one of the theorems 8 or 9.
\end{examinstructions}
\section{III.}
\begin{examinstructions}
Solve the following problems:
\end{examinstructions}
\begin{examinstructions}
\((X,\Sigma,\mu)\) is a measure space. \(E\) is a Banach space.
\end{examinstructions}
\begin{examproblems}{0}
\item
Let \((f_n)\) be a sequence of \(L_E^1(\mu)\) such that \(\int \lVert f_n\rVert\,d\mu < \frac{1}{n^2}\) for every~\(n\). Prove or disprove the pointwise convergence and the convergence in the mean of the series \(\sum f_n\).
\item
Assume~\(\mu\) is finite and let~\(R\) be a ring generating~\(\Sigma\). Define \(\rho(A,B)=\mu(A\mathbin{\triangle}B)\) for \(A,B\in\Sigma\). Prove that~\(\rho\) is a semi-distance and that~\(R\) is dense in~\(\Sigma\) for~\(\rho\).
\item
Let \(X=\mathbb{R}\) and~\(\mu\) the Lebesgue measure. Let \(f\in L_{\mathbb{C}}^1(\mu)\) such that \(\int_0^x f\,d\mu=0\) for every \(x\in\mathbb{R}\) (if \(x<0\), \(\int_0^x=-\int_x^0\)). Prove that \(f=0\), \(\mu\)-a.e.
\item
Let \(\mu,\nu\) be finite, \(\sigma\)-additive measures on~\(\Sigma\), such that \(\mu\ll\nu\) and \(\nu\ll\mu\). What can be said about the relationship of the Radon–Nikodym densities of one measure with respect to the other.
\item
Let \((x_n)\) be a sequence in the Banach space~\(E\). Assume \(x_n\to x\) weakly in~\(E\). Show that~\(x\) is the only weak limit of the sequence \((x_n)\) and that \(\sup_n\lVert x_n\rVert<\infty\). (Hint: Consider the \(x_n\) as elements of the bidual \(E^{**}\).)
\item
Let \(X=\mathbb{R}\) and~\(\mu\) the Lebesgue measure. Let \(f\in L^1(\mu)\) and for each \(t\in\mathbb{R}\) define \(f_t(x)=f(x-t)\). Show that the mapping \(t\mapsto f_t\) is a continuous map from~\(\mathbb{R}\) into \(L^1(\mu)\). (Hint: prove it first for continuous functions with compact support).
\end{examproblems}
\end{document}
